\documentclass[aps,prd,reprint,superscriptaddress,twocolumn]{revtex4-2}
\usepackage{graphicx}
\usepackage{amsmath,amssymb}
\usepackage{bm}
\usepackage{hyperref}

\begin{document}

\title{Correlation Quadrature Uncertainty as a Falsification Criterion for Standard Quantum Mechanics}

\author{Author Names}
\email{email@institution.edu}
\affiliation{Institution}

\begin{abstract}
We derive an exact inequality for fermion many-body systems: the product of correlation quadrature uncertainties is bounded below by the inter-site density gradient, $\Delta C^R \cdot \Delta C^I \geq \frac{1}{4}|\Delta n|$. This inequality follows from the three constitutive assumptions of standard quantum mechanics---Fermi-Dirac statistics, unitary evolution, and Gaussian state structure---and is therefore a necessary condition for their joint validity. Its experimental violation signals a breakdown of standard quantum mechanics at a precisely identifiable location. The inequality has been verified for 85,639 quantum states ($L = 2$--$8$; ground states and non-equilibrium steady states; 0\% violation rate), establishing a robust null hypothesis. We identify three experimentally accessible systems where violation may occur: fractional quantum Hall edge states ($\nu = 1/3$), where anyonic statistics generically modify the fundamental commutator; strongly driven fermion chains, where non-Gaussian correlations break Wick's theorem at the 10--30\% level; and BEC analogue black holes, where the $\mathcal{J} = 0$ information-conservation constraint produces non-diagonal frequency correlations $C_J(\omega,\omega') \neq 0$. The inequality is tight---equality is reached for pure states with real $\alpha\beta^*$. Any statistically significant violation constitutes a discovery of physics beyond standard quantum mechanics; a null result maps the validity boundary of the standard framework in previously untested regimes.
\end{abstract}

\maketitle

% Content will be filled from PRD_Letter_提交稿.md after review rounds complete

\end{document}
